% !TeX program = xelatex
\documentclass[11pt,a4paper]{article}
\usepackage[margin=23mm]{geometry}
\usepackage{fontspec}
\setmainfont{Latin Modern Roman}
\usepackage{amsmath,amssymb,mathtools}
\usepackage{xcolor,enumitem,booktabs,longtable,array,xurl,hyperref}
\definecolor{muted}{HTML}{53636C}
\hypersetup{colorlinks=true,urlcolor=blue,linkcolor=blue}
\setlength{\parindent}{0pt}
\setlength{\parskip}{5pt}
\newcommand{\R}{\mathbb R}
\newcommand{\E}{\mathbb E}
\newcommand{\N}{\mathbb N}
\newcommand{\ind}{\mathbf 1}
\newcommand{\Var}{\operatorname{Var}}
\newcommand{\Cov}{\operatorname{Cov}}
\newcommand{\supp}{\operatorname{supp}}
\DeclareMathOperator*{\argmin}{arg\,min}
\DeclareMathOperator*{\argmax}{arg\,max}
\newcommand{\entrymeta}[3]{{\sffamily\small\color{muted}\textbf{\texttt{#1}}\quad|\quad #2\par #3\par}}
\newcommand{\Problemref}[1]{\texttt{#1}}
\newcommand{\JELref}[2]{\texttt{#2} (JEL \texttt{#1})}
\begin{document}
\section*{Gross substitutes: truthfulness versus communication constraints}
\textbf{Problem ID:} \texttt{D82-94}\quad \textbf{JEL:} \texttt{D82}\par
% BEGIN ECONOMICS STATEMENT
\label{problem:OP-0174}
\label{game:GT-044}
{\sffamily\small\color{muted}
\textbf{GT-044} | Computational mechanism design | \textbf{H}: unrestricted catalogue claim is already settled; constrained status requires care.
\par}
\subsection*{Definitions and mathematical statement}
Let $m,n\geq1$ be integers, let $M$ be a finite set with $|M|=m$, and index bidders by $[n]=\{1,\ldots,n\}$. Each valuation $v_i:2^M\to\mathbb R_{\geq0}$ is normalized, meaning $v_i(\varnothing)=0$, and monotone, meaning $v_i(S)\leq v_i(T)$ whenever $S\subseteq T$. For $p\in\mathbb R^M$, define $D_i(p)=\operatorname*{arg\,max}_{S\subseteq M}(v_i(S)-\sum_{g\in S}p_g)$. Require gross substitutes: for every $p,p'\in\mathbb R^M$ with $p'\geq p$ coordinatewise and every $S\in D_i(p)$, some $S'\in D_i(p')$ contains $S\cap\{g:p'_g=p_g\}$. An allocation $A=(A_1,\ldots,A_n)$ consists of pairwise disjoint subsets of $M$; items may remain unallocated. Its welfare is $W(A,v)=\sum_{i=1}^n v_i(A_i)$, and $\operatorname{OPT}(v)=\max_A W(A,v)$. Both suprema below range over profiles of these normalized monotone gross-substitutes valuations; the bounded-communication domain also obeys the value-length restriction below. In both ratios, a positive numerator divided by zero achieved welfare is $+\infty$.
For unrestricted direct DSIC mechanisms, the original ratio is exactly
\[
 \inf_{\mathcal M\text{ DSIC}}\sup_{v:\operatorname{OPT}(v)>0}
 \frac{\operatorname{OPT}(v)}{\mathbb E W(\mathcal M(v),v)}=1.
\]
The expectation concerns internal mechanism randomization, if any. For this unrestricted benchmark, randomized direct DSIC means truthful reporting maximizes expected quasilinear utility against every profile of other reports; restricting instead to deterministic mechanisms or universally truthful mechanisms gives the same benchmark value.
The welfare-maximizing VCG mechanism establishes this equality. The source's substantive question instead concerns protocols with bounded communication. Given an integer communication budget $C\geq0$, define
\[
 \rho(m,n,C)=\inf_{(\Pi,s)}\sup_{v:\operatorname{OPT}(v)>0}
 \frac{\operatorname{OPT}(v)}{W(A^\Pi(s(v)),v)},
\]
where $\Pi$ is a deterministic finite public-blackboard protocol. Its rooted tree specifies at each internal node which bidders speak simultaneously, their legal binary messages and the next node for each joint message. Message boundaries are fixed by the protocol, with any information conveyed by message lengths charged to communication. Every message is revealed publicly after its round; a message cannot depend on another message in the same round. The total communicated bits on every legal root-to-leaf path, including paths generated by deviations, is at most $C$. Each leaf labels an allocation $A^\Pi$ and a real payment vector $P^\Pi$. Leaf labels and the protocol description are fixed independently of the private valuations; their description lengths are not charged as communication. Local computation is unrestricted.
For each admissible $v_i$, $s_i(v_i)$ is a pure strategy specifying a legal message at every node where bidder $i$ speaks. Writing $A^\Pi(t),P^\Pi(t)$ for the outcome of a pure strategy profile $t$, require, for every bidder $i$, every admissible $v_i$, every legal pure opponent profile $t_{-i}$ and every legal pure deviation $t_i$,
\[
 v_i(A_i^\Pi(s_i(v_i),t_{-i}))-P_i^\Pi(s_i(v_i),t_{-i})
 \geq v_i(A_i^\Pi(t_i,t_{-i}))-P_i^\Pi(t_i,t_{-i}).
\]
Thus $s_i(v_i)$ is dominant against all contingent opponent strategies. The pure-strategy inequalities also imply the corresponding expected-utility inequalities for randomized strategies in this finite protocol. Welfare in $\rho$ is evaluated at the prescribed profile $s(v)$, without requiring the same welfare at every dominant-strategy equilibrium. Require $P_i^\Pi=0$ whenever $A_i^\Pi=\varnothing$ at every leaf. No individual-rationality, nonnegative-payment or participation constraint is imposed beyond this normalization. Values have at most $m^k$ bits per bundle for a fixed integer $k\geq1$; dependence on this fixed $k$ and on the value encoding is suppressed in $\rho$. This per-bundle restriction does not assume a succinct representation of the entire valuation. Determine the tradeoff as $C$ varies, particularly for families of nonnegative integer budgets $C(m,n)$ bounded by a polynomial in $m,n$. Bounds should state their parameter ranges and any dependence on $k$, the value encoding and the chosen budget polynomial; the question imposes no running-time bound.
\subsection*{Short English statement}
The unrestricted optimum is known. The meaningful question is how much welfare dominant-strategy communication protocols can achieve with limited messages.
\subsection*{Variants and scope boundaries}
Low-communication ex-post implementations, direct DSIC with a demand oracle, universally truthful randomized protocols, and truthfulness in expectation have different strategic constraints. They cannot be substituted into the same infimum.
\subsection*{Source relationship and status}
[S1] explicitly poses the constrained approximation question. An indexed October 2026 preprint [S2] reports new GS lower bounds answering part of it; its full text could not be verified here. Consequently the constrained question is recorded at its 2022 source date, with subsequent status unverified, while the original unrestricted claim is corrected outright.
The preceding verification qualification and those in the references record the original source assessment. In a subsequent full-text consultation, [S2], Theorem 1.1 and Sections 2.1--2.2, state a GS lower bound under individual rationality, no negative transfers and a different succinct value encoding. Those hypotheses are not imposed in $\rho$ here, and no transfer of that bound or complete characterization of $\rho(m,n,C)$ is asserted. This comparison checks the reported formulation, not its proof.
\subsection*{References}
\begingroup\footnotesize\raggedright
\textbf{[S1]} Shahar Dobzinski, Shiri Ron, and Jan Vondr\'ak (2022). \emph{On the Hardness of Dominant Strategy Mechanism Design}. arXiv:2206.00334v1. \textit{Verified locator:} introduction, ``Open Questions and Future Directions,'' printed pp. 4--5; Section 8, Definition 8.1 and Theorem 3.2. \url{https://arxiv.org/pdf/2206.00334}
\par\textbf{[S2]} \emph{The Hardness of Dominant Strategy Mechanism Design, Revisited} (2026), arXiv:2610.06773. \textit{Verification limit:} indexed abstract only; reported GS resolution must be checked in the full manuscript before publication. \url{https://arxiv.org/abs/2610.06773}
\par\endgroup

\subsection*{Common conventions: Source mathematical conventions}

\textbf{Finite games.} For a finite set $A$, $\Delta(A)$ is its probability
simplex. Players choose independent mixed strategies in Nash equilibrium;
correlation is allowed only when explicitly part of the solution concept.
In a game with payoff functions $u_i$ and strategy profile $x$, an additive
$\varepsilon$ Nash equilibrium satisfies
\[
 u_i(x)\geq u_i(x_i',x_{-i})-\varepsilon
 \quad\text{for every player $i$ and every admissible deviation $x_i'$}.
\]
For numerical approximation, payoffs are normalized as locally specified.
An $\varepsilon$ payoff residual is distinct from distance at most
$\varepsilon$ to the set of exact equilibria. Point convergence, convergence
of empirical play, and convergence in distance to an equilibrium set are
also distinct.

\textbf{Computational representations.} Unless stated otherwise, a complexity
question takes rational data with binary encoding; polynomial time is measured
in total bit length. A fixed-accuracy polynomial algorithm is not necessarily
polynomial in $1/\varepsilon$, and neither is necessarily polynomial in
$\log(1/\varepsilon)$. Succinct games and value, demand, or sampling oracles
must declare their interfaces. Exact real solutions require an explicit
representation; an unspecified list of arbitrary real numbers is not a
finite computational output. A claimed algorithmic barrier is meaningful
only for its specified model and reductions.

\textbf{Dynamic games.} Histories, information, observation, and allowable
randomization are part of the model. A stationary strategy depends only on
the current state; a positional strategy in a deterministic graph game
chooses one action at each controlled vertex. Nash equilibrium at an initial
state and equilibrium after every history are different requirements.
An expectation of a pathwise $\liminf$ payoff, a $\liminf$ of expected averages,
a limiting expected average, and a uniform finite-horizon equilibrium are
different criteria. None is silently substituted for another.

\textbf{Allocation and mechanisms.} A complete allocation assigns every item
exactly once unless otherwise stated. Zero-valued goods, negative values,
capacity constraints, feasibility, lotteries, and copies of goods affect
fairness definitions and are specified locally. Dominant-strategy incentive
compatibility, Bayesian incentive compatibility, universal truthfulness,
and truthfulness in expectation impose different inequalities. Whenever
an objective ratio has zero denominator, its convention is stated or those
instances are excluded.

\textbf{Infinite-dimensional models.} Expectations, integrals, stochastic
equations, and optimization objectives require their local regularity,
integrability, filtrations, and admissible domains. A backward--forward
mean-field system is a boundary-value problem; it is not an initial-value
semiflow without an additional construction. Long-time, large-population,
and vanishing-noise limits require an order of limits and a convergence
topology. If a minimum need not be attained, the objective is its infimum
or supremum and the approximation requirement is stated separately.
% END ECONOMICS STATEMENT
\end{document}
